\documentclass[leqno]{article}
% Shared QBT/Soln pack preamble for resources.danieljsmith.org.
\usepackage{graphicx}
\usepackage{dsfont}
\usepackage{mathtools}
\usepackage{amssymb}
\usepackage{amsmath}
\usepackage{mleftright}
\usepackage{amsthm}
\usepackage{float}
\usepackage{ragged2e}
\usepackage{tensor}
\usepackage{fancyhdr}
\usepackage{empheq}
\usepackage[most]{tcolorbox}
\usepackage{enumitem}
\usepackage{dirtytalk}
\usepackage{marginnote}
\usepackage{microtype}
\usepackage[
  a4paper,
  left=0.8in,
  right=1in,
  top=1in,
  bottom=0.5in,
  headheight=14pt,
  headsep=0.4in
]{geometry}
\setlength{\parindent}{0cm}

% TikZ / PGFPlots
\usepackage{pgfplots}
\pgfplotsset{compat=1.18}
\usepgfplotslibrary{fillbetween, polar}
\usetikzlibrary{calc, positioning}
\usetikzlibrary{angles, quotes}
\usetikzlibrary{arrows.meta}
\usetikzlibrary{decorations.pathreplacing, decorations.markings}
\usetikzlibrary{patterns}
\usetikzlibrary{intersections}

% Links and contents. The contents follow the TMUA papers: compact entries,
% no dot leaders, and each question's mark total in grey after its label.
\usepackage{tocloft}
\tocloftpagestyle{fancy}
\setlength{\cftbeforesubsecskip}{2pt}
\setlength{\cftsubsecindent}{1.5em}
\setlength{\cftsubsecnumwidth}{0pt}
\renewcommand{\cftsubsecfont}{\small\color{black}}
\renewcommand{\cftsubsecpagefont}{\small\color{black}}
\renewcommand{\cftsubsecleader}{\hfill}
\usepackage[colorlinks=true,linkcolor=black,citecolor=magenta,urlcolor=black]{hyperref}%
\urlstyle{same}

% Pack macros
% \marks also adds its value to the current question's total in the .aux
% file, so the contents show each total from the second pass. The write
% happens at shipout, so a box typeset more than once counts once. The
% pack's mark total for the contents-page title is summed the same way.
\makeatletter
\newcount\djs@totalmarks
\newcommand{\djsAddMarks}[2]{%
  \global\advance\djs@totalmarks#2\relax
  \@ifundefined{djs@marks@#1}%
    {\expandafter\xdef\csname djs@marks@#1\endcsname{#2}}%
    {\expandafter\xdef\csname djs@marks@#1\endcsname
      {\the\numexpr\csname djs@marks@#1\endcsname+#2\relax}}}
\RenewDocumentCommand{\marks}{o m}{%
  \protected@write\@auxout{}{\string\djsAddMarks{\number\value{enumi}}{#2}}%
  \IfValueTF{#1}%
    {\marginnote{\raggedright\textbf{[#2]}}[#1]}%
    {\marginnote{\raggedright\textbf{[#2]}}}%
}
\makeatother

\newcommand{\vspacerule}[2]{%
  \par\vspace*{#1}%
  \noindent\hrulefill\par
  \vspace*{#2}%
}

\definecolor{scarlet}{RGB}{205, 40, 75}
% Site colours for the contents-page title: the accent, the rule colour, a
% stronger rule colour for the path separator, and the muted text colour.
\definecolor{djsBurgundy}{RGB}{118, 59, 67}
\definecolor{djsHairline}{RGB}{216, 213, 205}
\definecolor{djsSlash}{RGB}{170, 166, 158}
\definecolor{djsMuted}{RGB}{96, 101, 106}

% Worked-solution box, adapted from the TMUA worked-solutions box. A white breakable box with a left scarlet rule and no surrounding
% frame or tint. A part that breaks fills the rest of its page, so the rule
% runs to the bottom margin; the unbroken or last part ends in a short
% horizontal foot rule so a solution visibly stops. Unlike the TMUA box there
% is no answer label. The box does not grow to the left, so the rule sits at
% the text edge; growing to the right keeps the usable line width close to the
% old tcolorbox Solution box.
% Spacing inside the box: paragraphs are 6pt apart, so source blank lines show
% (solutions put one after each display and none between a lead-in and its
% display). Every display gets 6pt above and 10pt below, whatever the length
% of its lead-in line; align rows are spread by an 8pt \jot; solution parts
% are 14pt apart and sub-parts 8pt; \left and \right space like ordinary
% brackets (mleftright); no single line is stranded at a page break; and
% neither the Solution label nor the line before a list ends a page. No
% skip can shrink, so full pages are not squeezed.
\newtcolorbox{djsSolution}{%
  enhanced, breakable, frame hidden, sharp corners, colback=white,
  height fixed for=first and middle, height fill,
  extras unbroken and last={natural height},
  borderline west={2.2pt}{0pt}{scarlet},
  left=11pt, right=0pt, top=5pt, bottom=13pt,
  before skip=13pt, after skip=4pt,
  grow to right by=10mm,
  overlay unbroken and last={%
    \fill[scarlet] (frame.south west)
      rectangle ([xshift=3.5cm, yshift=2.2pt]frame.south west);},
  before upper={{\color{scarlet}\bfseries Solution}\par\nobreak\vspace{4pt}%
    \setlength{\parskip}{6pt plus 2pt}%
    \setlength{\abovedisplayskip}{6pt plus 2pt}%
    \setlength{\abovedisplayshortskip}{6pt plus 2pt}%
    \setlength{\belowdisplayskip}{10pt plus 2pt}%
    \setlength{\belowdisplayshortskip}{10pt plus 2pt}%
    \setlength{\jot}{8pt}%
    \setlist{parsep=6pt plus 2pt, topsep=6pt plus 2pt, beginpenalty=10000}%
    \setlist[enumerate,1]{itemsep=14pt plus 2pt}%
    \setlist[enumerate,2]{itemsep=8pt plus 2pt}%
    \mleftright
    \clubpenalty=10000 \widowpenalty=10000 \displaywidowpenalty=10000\relax}}

% Optional smaller figure inside a solution box. \scalebox shrinks node text
% with the drawing. Not applied to existing figures.
\newcommand{\SolnFigureScale}{0.85}
\newsavebox{\djsSolnFigureBox}
\newenvironment{djsSolutionFigure}
  {\par\vspace{4pt}\begin{center}\begin{lrbox}{\djsSolnFigureBox}}
  {\end{lrbox}\scalebox{\SolnFigureScale}{\usebox{\djsSolnFigureBox}}%
   \end{center}\vspace{2pt}\par}

\newcommand{\cxl}[2]{%
  \tikz[baseline=(X.base)]{%
    \node[inner sep=0pt, outer sep=0pt, text=#1] (X) {$#2$};
    \draw[#1, line width=0.4pt] (X.south west)--(X.north east);
  }%
}

% Page-style hook run at the contents page and each question's opening page.
% A no-op for QBT packs; \djsSolnHeader points it at the fancy header.
\newcommand{\djsQuestionPageStyle}{}

% A contents entry: "Question N" in a box wide enough for two digits, so the
% mark totals line up, then the total in grey. Bookmarks keep "Question N".
\makeatletter
\newlength{\djsTocLabelWidth}
\newcommand{\djsTocMarks}[1]{%
  \@ifundefined{djs@marks@#1}{}{%
    {\footnotesize\itshape\color{black!45}%
     \csname djs@marks@#1\endcsname~mark%
     \ifnum\csname djs@marks@#1\endcsname=1 \else s\fi}}}
\newcommand{\djsTocEntry}[1]{%
  \texorpdfstring
    {\settowidth{\djsTocLabelWidth}{Question 00}%
     \makebox[\dimexpr\djsTocLabelWidth+0.9em\relax][l]{Question #1}%
     \djsTocMarks{#1}}%
    {Question #1}}
\makeatother

% The question count for the contents-page title is written to the .aux at
% the end of the document, so it also appears from the second pass.
\newcounter{djsQuestions}
\newcommand{\questionitem}{%
  \djsQuestionPageStyle
  \item\phantomsection
  \stepcounter{djsQuestions}%
  \addcontentsline{toc}{subsection}{\protect\djsTocEntry{\number\value{enumi}}}%
}
\makeatletter
\newcommand{\djsSetQuestionCount}[1]{\gdef\djs@questioncount{#1}}
\AtEndDocument{\immediate\write\@auxout
  {\string\djsSetQuestionCount{\number\value{djsQuestions}}}}
\makeatother

\renewcommand{\contentsname}{Questions}

% \djsPackHeader[<link content tag>]{<left header>}; the tag says which
% document a click on the header link came from.
\newcommand{\djsPackHeader}[2][questions]{%
  \pagestyle{fancy}%
  \fancyhead{}%
  \fancyfoot{}%
  \renewcommand{\headrulewidth}{0.9pt}%
  \lhead{#2}%
  \chead{}%
  \rhead{\href{https://danieljsmith.org/?utm_source=pdf&utm_campaign=further-maths&utm_content=#1}{danieljsmith.org}}%
}

% The header's topic is also the contents-page title, under which a label
% says which booklet this is.
\newcommand{\djsPackTitle}{}
\newcommand{\djsBookletLabel}{}
\newcommand{\djsQbtHeader}[1]{%
  \djsPackHeader{#1}%
  \renewcommand{\djsPackTitle}{#1}%
  \renewcommand{\djsBookletLabel}{Questions by topic}%
}

% The colophon at the foot of page one, independent of the page style: the
% contents-page title's hairline led by a short scarlet rule, then the date on
% the left and a link to the site on the right, all inside the bottom margin.
% Its sizes are fixed, so it matches the TMUA papers' colophon. The document
% supplies the date; recompiling never changes it.
\makeatletter
\newcommand{\djsLastUpdated}[1]{%
  \AtBeginDocument{\AddToHookNext{shipout/foreground}{\djs@colophon{#1}}}}
\newcommand{\djs@colophon}[1]{%
  \put(0.8in,-\dimexpr\paperheight-0.47in\relax){%
    \makebox[0pt][l]{\color{djsHairline}%
      \rule[0.9pt]{\dimexpr\paperwidth-1.8in\relax}{0.6pt}}%
    \makebox[0pt][l]{\color{scarlet}\rule{3.5cm}{2.2pt}}}%
  \put(0.8in,-\dimexpr\paperheight-0.32in\relax){\makebox[0pt][l]{%
    \fontsize{9}{11}\selectfont Last updated #1}}%
  \put(\dimexpr\paperwidth-1in\relax,-\dimexpr\paperheight-0.32in\relax){%
    \makebox[0pt][r]{\fontsize{7}{8.5}\selectfont\color{djsMuted}%
      Latest version at \href{https://resources.danieljsmith.org/}%
        {\color{black}resources.danieljsmith.org}}}}
\makeatother

% Solutions packs show the header only on the contents page and each
% question's opening page, so it never appears part-way through a solution.
% \thispagestyle marks the next page shipped out, which is the question's
% page because every solution question starts after a \newpage.
\newcommand{\djsSolnHeader}[1]{%
  \djsPackHeader[solutions]{\textit{Solutions} - #1}%
  \pagestyle{empty}%
  \renewcommand{\djsQuestionPageStyle}{\thispagestyle{fancy}}%
  \renewcommand{\djsPackTitle}{#1}%
  \renewcommand{\djsBookletLabel}{Worked solutions}%
}

% \djsSitePath{<strand>}{<topic>}, called once after the header, names where
% the pack sits on the site; the contents-page title prints it above the pack
% title, e.g. CORE PURE / FURTHER CALCULUS. Without it the line is left out.
\makeatletter
\newif\ifdjs@sitepath
\newcommand{\djs@strand}{}
\newcommand{\djs@section}{}
\newcommand{\djsSitePath}[2]{%
  \djs@sitepathtrue
  \renewcommand{\djs@strand}{#1}%
  \renewcommand{\djs@section}{#2}%
}
\makeatother

% Further Mechanics packs call \djsMechanicsPack in their preamble, so that
% \djsFrontMatter prints the default value of g above the contents. The box
% keeps its own space above it (before skip is the whole gap below the title,
% not an addition to it). Without the box, \djsGravityNote is just extra space
% before the contents heading, bringing the gap above it close to the gap
% between the box and the heading on Further Mechanics pages.
\newcommand{\djsGravityNote}{\vspace{14pt}}
\newcommand{\djsMechanicsPack}{%
  \renewcommand{\djsGravityNote}{%
    \begin{tcolorbox}[sharp corners, colback=white, colframe=black,
      boxrule=1pt, left=8pt, right=8pt, top=6pt, bottom=6pt,
      before skip=20pt]
    {\bfseries Instructions}\par\vspace{4pt}
    Unless a question states otherwise, take $g = 9.8\,\mathrm{m\,s^{-2}}$
    wherever a numerical value of $g$ is required. Solutions also give final
    answers for $g = 10\,\mathrm{m\,s^{-2}}$ where these differ.
    \end{tcolorbox}%
    \vspace{16pt}}%
}

% Contents-page title: the site path in spaced burgundy capitals, the pack
% title, a hairline led by a short scarlet rule (the solution box's foot rule),
% then the booklet label with the question and mark counts. The counts come
% from the .aux, so they appear from the second pass.
\makeatletter
\newsavebox{\djs@eyebrowbox}
\newcommand{\djs@caps}[1]{\textls[180]{\MakeUppercase{#1}}}
\newcommand{\djs@eyebrow}{%
  \color{djsBurgundy}\footnotesize\bfseries
  \djs@caps{\djs@strand}%
  {\color{djsSlash}\hspace{0.55em}/\hspace{0.55em}}%
  \djs@caps{\djs@section}}
\newcommand{\djs@dot}{\hspace{0.45em}\textperiodcentered\hspace{0.45em}}
\newcommand{\djs@counts}{%
  \@ifundefined{djs@questioncount}{}{%
    \djs@questioncount~question\ifnum\djs@questioncount=1 \else s\fi
    \ifnum\djs@totalmarks>0
      \djs@dot\the\djs@totalmarks~mark\ifnum\djs@totalmarks=1 \else s\fi
    \fi}}
\newcommand{\djsContentsTitle}{%
  \ifdjs@sitepath
    \sbox\djs@eyebrowbox{\djs@eyebrow}%
    \ifdim\wd\djs@eyebrowbox>\linewidth
      \@latex@warning@no@line{The \string\djsSitePath\space line is wider
        than the text, so it wraps. Shorten its topic}%
    \fi
    {\raggedright\djs@eyebrow\par}%
    \vspace{7pt}%
  \fi
  {\raggedright\fontsize{24.88}{28}\selectfont\djsPackTitle\par}%
  \vspace{9pt}%
  \noindent\makebox[0pt][l]{\color{djsHairline}\rule[0.9pt]{\linewidth}{0.6pt}}%
  {\color{scarlet}\rule{3.5cm}{2.2pt}}\par
  \vspace{5pt}%
  {\small\color{djsMuted}\djsBookletLabel\hfill\djs@counts\par}%
  \vspace{8pt}%
}
\makeatother

\newcommand{\djsFrontMatter}{%
  \djsQuestionPageStyle
  \vspace*{6pt}%
  \djsContentsTitle
  \djsGravityNote
  \tableofcontents
  \newpage
}


\djsQbtHeader{Impulse and Work Done by a Variable Force}
\djsSitePath{Further Mechanics}{Kinematics \& Dynamics}
\djsMechanicsPack
\djsLastUpdated{15 October 2026}

\begin{document}
\djsFrontMatter
\begin{enumerate}[
  label=\textbf{\arabic*.},
  leftmargin=*,
  topsep=0pt,
  partopsep=0pt,
  parsep=0pt
]

% ---- Question 1 ----
\questionitem
% [Source: _QBT__Impulse_and_Work_Done_by_a_Variable_Force, Q1]

A puck of mass $0.40\,\mathrm{kg}$ moves on a smooth horizontal surface. \\
The horizontal unit vectors $\mathbf{i}$ and $\mathbf{j}$ are perpendicular. \\

A hockey stick strikes the puck. \\
Immediately before contact, the puck has velocity $9\mathbf{i}\,\mathrm{m\,s^{-1}}$. \\
Immediately after contact, its velocity is $12\mathbf{j}\,\mathrm{m\,s^{-1}}$. \\
\begin{enumerate}[label=\textbf{(\alph*)}]
\item Find the impulse exerted by the stick on the puck, giving your answer as a vector, and calculate its magnitude. \marks{4} \\
\item During contact, the force exerted by the stick acts in a constant horizontal direction. \\
Its magnitude, $F$ newtons, is modelled by
\[
F=k\sin(25\pi t),\qquad 0\leq t\leq 0.040
\] \\
where $k$ is a positive constant and $t$ is the time in seconds measured from the start of contact. \\
The stick loses contact with the puck at $t=0.040$. \\
Using your answer to part (a), find the exact value of $k$. \marks{3} \\
\end{enumerate}

\newpage

% ---- Question 2 ----
\questionitem
% [Source: _QBT__Impulse_and_Work_Done_by_a_Variable_Force, Q2]

A trolley of mass $6\,\mathrm{kg}$ moves along a smooth horizontal track at $4\,\mathrm{m\,s^{-1}}$ and strikes a buffer. \\
The buffer brings the trolley to rest without causing it to rebound. \\
Air resistance is negligible. \\
\begin{enumerate}[label=\textbf{(\alph*)}]
\item Find the magnitude of the impulse exerted by the buffer on the trolley. \marks{2} \\
\item The trolley first touches the buffer at time $t=0$ and comes to rest at time $t=T$ seconds. \\
During this interval, the buffer exerts a horizontal force opposite to the trolley's motion. \\
The magnitude of this force is
\[
F(t)=18\,000t(T-t)\,\mathrm{N},\qquad 0\leq t\leq T
\] \\
Using your answer to part (a), find the stopping time $T$ and the maximum magnitude of the force exerted by the buffer. \marks{5} \\
\end{enumerate}

\newpage

% ---- Question 3 ----
\questionitem
% [Source: _QBT__Impulse_and_Work_Done_by_a_Variable_Force, Q3]

A sled of mass $20\,\mathrm{kg}$ slides down a smooth straight slope inclined at $30^\circ$ to the horizontal. \\
It reaches a hydraulic buffer with speed $3\,\mathrm{m\,s^{-1}}$. \\

The buffer exerts a force on the sled directed up the slope. \\
When the buffer has been compressed by $x$ metres, the magnitude of this force is $F$ newtons, where
\[
F=\frac{A}{1+2x}+200x
\] \\
and $A$ is a positive constant. \\
The sled comes to rest after compressing the buffer by $0.5\,\mathrm{m}$. \\
Throughout the compression, the sled moves along the slope. \\
Model the sled as a particle. \\

\begin{enumerate}[label=\textbf{(\alph*)}]
  \item Find, in terms of $A$, the work done against the buffer during this compression. \marks{3} \\
  \item Using work and energy, find the exact value of $A$, taking account of the work done by gravity. \marks{3} \\
\end{enumerate}

\newpage

% ---- Question 4 ----
\questionitem
% [Source: _QBT__Impulse_and_Work_Done_by_a_Variable_Force, Q4]

A trolley of mass $80\,\mathrm{kg}$ travels along a straight horizontal track. \\
A powered launcher exerts a variable driving force on the trolley in the direction of motion for $2\,\mathrm{s}$. \\
Throughout this interval, a constant resistance of $40\,\mathrm{N}$ acts against the motion. \\

At time $t=0$, when the launcher begins to act, the trolley has speed $3\,\mathrm{m\,s^{-1}}$. \\
At time $t=2\,\mathrm{s}$, its speed is $7\,\mathrm{m\,s^{-1}}$. \\
The trolley does not change its direction of motion. \\
\begin{enumerate}[label=\textbf{(\alph*)}]
\item Find the magnitude of the impulse exerted on the trolley by the launcher during this interval. \marks{3} \\
\item The driving force exerted by the launcher is modelled by
\[
F=kt^2(2-t),\qquad 0\leq t\leq 2
\] \\
where $F$ is measured in newtons, $t$ is measured in seconds and $k$ is a positive constant. \\
\begin{enumerate}[label=(\roman*)]
\item Find the value of $k$. \marks{3} \\
\item Determine the maximum driving force predicted by this model. \marks{2} \\
\end{enumerate}
\end{enumerate}

\newpage

% ---- Question 5 ----
\questionitem
% [Source: _QBT__Impulse_and_Work_Done_by_a_Variable_Force, Q5]

A particle $P$ of mass $2\,\mathrm{kg}$ moves along a straight horizontal line. \\
Take the positive direction along the line to be the direction of its initial motion. \\

At time $t=0$, the velocity of $P$ is $3\,\mathrm{m\,s^{-1}}$. \\
For $0\leq t\leq 2$, the only horizontal force acting on $P$ has signed component
\[
F(t)=(3t^2-12t+9)\,\mathrm{N}
\] \\
in the positive direction. \\
The particle continues to move in the positive direction throughout this interval. \\
\begin{enumerate}[label=\textbf{(\alph*)}]
\item Find the impulse of the force over the interval $0\leq t\leq 2$. \marks{3} \\
\item Find the speed of $P$ at $t=2$. \marks{2} \\
\item Find the work done by the force during the portion of the interval $0\leq t\leq 2$ for which the force opposes the motion of $P$. \marks{5} \\
\end{enumerate}

\newpage

% ---- Question 6 ----
\questionitem
% [Source: _QBT__Impulse_and_Work_Done_by_a_Variable_Force, Q6]

A particle $P$ of mass $4\,\mathrm{kg}$ moves along a straight line on a smooth horizontal surface. \\
At time $t=0$, its velocity is $4\,\mathrm{m\,s^{-1}}$ in the positive direction. \\
The only horizontal force acting on $P$ has component
\[
F(t)=(3t^2-12t)\,\mathrm{N}
\] \\
in the positive direction, where $t$ is the time in seconds and $t\geq 0$. \\
\begin{enumerate}[label=\textbf{(\alph*)}]
\item Find the magnitude of the impulse exerted by the force between $t=0$ and $t=2$. \marks{3} \\
\item Find the speed of $P$ at $t=3$, and state whether it is moving in the positive or negative direction. \marks{3} \\
\item Find the exact time at which $P$ next comes to rest after $t=2$. \marks{4} \\
\end{enumerate}

\newpage

% ---- Question 7 ----
\questionitem
% [Source: _QBT__Impulse_and_Work_Done_by_a_Variable_Force, Q7]

A particle of mass $4\,\mathrm{kg}$ moves along a straight horizontal line. \\
A fixed positive direction is chosen along the line. \\
At time $t$ seconds, where $t\geq 0$, the component of the resultant force in this direction is
\[
F(t)=6(t-1)(t-3)\,\mathrm{N}
\] \\
At time $t=0$, the velocity of the particle is $-1\,\mathrm{m\,s^{-1}}$. \\
\begin{enumerate}[label=\textbf{(\alph*)}]
\item Find the magnitude of the impulse exerted on the particle between $t=0$ and $t=1$. \marks{3} \\
\item Find the velocity of the particle at $t=1$. \marks{2} \\
\item Hence find all times $t>0$ at which the speed of the particle is $1\,\mathrm{m\,s^{-1}}$. \marks{6} \\
\end{enumerate}

\newpage

% ---- Question 8 ----
\questionitem
% [Source: _QBT__Impulse_and_Work_Done_by_a_Variable_Force, Q8]

In this question use $g=9.8\,\mathrm{m\,s^{-2}}$. \\

A compact load of mass $60\,\mathrm{kg}$ is raised vertically by a cable. \\
When the load is $x$ metres above its initial position, the tension in the cable is $P$ newtons, where
\[
P=498+\frac{k}{(x+2)^2}
\] \\
The constant $k$ satisfies $k>360$. \\
The load starts from rest at $x=0$, and the cable remains taut throughout the motion considered. \\
Air resistance is ignored. \\
\begin{enumerate}[label=\textbf{(\alph*)}]
\item Find, in terms of $k$ and $h$, the work done by the tension as the load rises from $x=0$ to $x=h$. Also find, in terms of $k$, the height at which the tension equals the weight of the load. \marks{5} \\
\item The load has speed $3\,\mathrm{m\,s^{-1}}$ when $x=6$. Determine $k$ and find the speed of the load when $x=8$, giving the speed to three significant figures. \marks{5} \\
\item Find the height at which the load first comes instantaneously to rest after leaving its initial position. \marks{2} \\
\item Give one physical reason why neglecting air resistance is reasonable for this load during its ascent. \marks{1} \\
\end{enumerate}

\newpage

% ---- Question 9 ----
\questionitem
% [Source: _QBT__Impulse_and_Work_Done_by_a_Variable_Force, Q9]

\begin{figure}[H]
    \centering
\begin{tikzpicture}[scale=1.2]
\def\ang{60}
\def\raylen{3.1}
\def\wallup{2.7}
\def\walldown{2.0}
\def\ballrad{0.1}
\def\arcRad{0.75}
\coordinate (P) at (0,0);
\coordinate (W) at ($(P)+(0,\wallup)$);
\coordinate (L) at ($(P)+(0,-\walldown)$);
\coordinate (B) at ($(P)+({90-\ang}:\raylen)$);
\coordinate (Ustart) at ($(P)!0.8!(B)$);
\coordinate (Uend) at ($(P)!0.5!(B)$);
\draw[thick, black] (L) -- (W);
\draw[thick, black] (P) -- (B);
\draw[-{Stealth[length=3mm]}, thick, black] (Ustart) -- (Uend) node[midway, below right] {$5\,\mathrm{m\,s^{-1}}$};
\draw[thick, black, fill=gray!30] (B) circle (\ballrad);
\pic["$60^\circ$", draw, angle radius=\arcRad cm, angle eccentricity=1.3] {angle=B--P--W};
\end{tikzpicture}
\end{figure}


A ball of mass $0.4\,\mathrm{kg}$ moves on a smooth horizontal surface and collides with a smooth fixed vertical wall. \\
Immediately before impact, its speed is $5\,\mathrm{m\,s^{-1}}$ and its direction makes an angle of $60^\circ$ with the wall, as shown in the diagram above. \\

The ball first touches the wall at time $t=0$ and loses contact at time $t=0.2\,\mathrm{s}$. \\
During contact, the magnitude of the force exerted by the wall on the ball is modelled by
\[
F=\begin{cases}
kt(0.2-t), & 0\leq t\leq 0.1\\
k(0.2-t)^2, & 0.1<t\leq 0.2
\end{cases}
\] \\
where $F$ is measured in newtons and $k$ is a positive constant. \\
\begin{enumerate}[label=\textbf{(\alph*)}]
\item Show, by integration, that the magnitude of the impulse exerted by the wall during the collision is
\[
\frac{k}{1000}\,\mathrm{N\,s}
\] \marks[-30pt]{5}
\item The coefficient of restitution between the ball and the wall is known to lie between $\frac{1}{4}$ and $\frac{3}{4}$, inclusive. \\
Determine the corresponding range of values of $k$, giving exact bounds and fully justifying your answer. \marks{5} \\
\item Given that $k=1500\sqrt{3}$, find the speed of the ball immediately after the collision in exact form, and the angle its direction of motion makes with the wall, to the nearest $0.1^\circ$. \marks{5} \\
\end{enumerate}

\newpage

% ---- Question 10 ----
\questionitem
% [Source: _QBT__Impulse_and_Work_Done_by_a_Variable_Force, Q10]

A car of mass $800\,\mathrm{kg}$ moves along a straight horizontal road. \\
At time $t=0$, it passes a marker with speed $20\,\mathrm{m\,s^{-1}}$. \\
After travelling a distance $x\,\mathrm{m}$ beyond the marker, its speed is $v\,\mathrm{m\,s^{-1}}$. \\

The driving force varies with distance and has magnitude
\[
\left(1800+1200e^{-x/200}\right)\,\mathrm{N}
\] \\
The resistance to motion has magnitude $2v^2\,\mathrm{N}$. \\

\begin{enumerate}[label=\textbf{(\alph*)}]
  \item Form a differential equation relating $v$ and $x$. \\
  By making the substitution $u=v^2$, solve this equation to show that
  \[
  v^2=900+(3x-500)e^{-x/200}
  \] \marks[-20pt]{8}
  \item Using the work--energy principle and your result from part (a), find the exact work done against resistance during the first $200\,\mathrm{m}$ beyond the marker. \\
  Hence find the average magnitude of the resistance over this distance. \marks{8} \\
\end{enumerate}

\newpage

% ---- Question 11 ----
\questionitem
% [Source: _QBT__Impulse_and_Work_Done_by_a_Variable_Force, Q11]

A particle of mass $2\,\mathrm{kg}$ moves along a horizontal straight line. \\
At time $t=0$, it passes through a point $O$ with velocity $1\,\mathrm{m\,s^{-1}}$. \\
The resultant force acts in the direction of motion and has magnitude $F\,\mathrm{N}$, where
\[
F=\frac{27}{v}-3v
\] \\
and $v\,\mathrm{m\,s^{-1}}$ is the velocity of the particle. \\
Its displacement from $O$ at time $t$ seconds is $x\,\mathrm{m}$. \\

\begin{enumerate}[label=\textbf{(\alph*)}]
  \item Formulate and solve a differential equation relating $v$ and $x$ to show that
  \[
  x=\ln\left(\frac{3+v}{2(3-v)}\right)-\frac{2}{3}(v-1),\qquad 1\leq v<3
  \] \marks[-40pt]{6}
  \item Find the exact displacement of the particle when its velocity first reaches $2\,\mathrm{m\,s^{-1}}$. \\
  By changing the variable of integration from $x$ to $v$, calculate the work done by the force up to this instant. \\
  Hence find the average force over this displacement, giving your answer to three significant figures. \marks{5} \\
  \item Formulate and solve a differential equation relating $v$ and $t$ to obtain an expression for $v$ in terms of $t$. \marks{5} \\
  \item Calculate the velocities at $t=1$ and $t=2$, giving your answers to four significant figures. \\
  Find the impulse of the force over the interval $1\leq t\leq2$, giving your answer to three significant figures. \marks{4} \\
\end{enumerate}

\end{enumerate}
\end{document}
